\documentclass[11pt]{article}

% Shared notes settings for Elements of AIML lectures (UPES).
% Keep this file free of \documentclass to allow reuse via \input.

\usepackage[utf8]{inputenc}
\usepackage[T1]{fontenc}
\usepackage[a4paper,margin=1in]{geometry}
\usepackage{hyperref}
\usepackage{xurl}
\usepackage{booktabs}
\usepackage{amsmath}
\usepackage{amssymb}
\usepackage{listings}
\usepackage{xcolor}
\usepackage{enumitem}
\usepackage{parskip}

% Code listings (general-purpose).
\lstset{
  basicstyle=\ttfamily\small,
  keywordstyle=\color{blue},
  commentstyle=\color{gray},
  stringstyle=\color{teal},
  showstringspaces=false,
  columns=fullflexible,
  frame=single,
  framerule=0.2pt,
  breaklines=true
}

\setlist[itemize]{topsep=0.3em,itemsep=0.2em}
\setlist[enumerate]{topsep=0.3em,itemsep=0.2em}

% Simple per-section source line.
\newcommand{\notesources}[1]{\par\noindent{\small\color{gray}\textit{Sources:} \sloppy #1\par}}

% Unit-wise source strings for this subject (used by slides/notes).
% Path is relative to the lecture `latex/` folder (the compilation CWD).
\IfFileExists{../../../_shared/aiml-sources.tex}{%
  % Source strings for Elements of AIML (used by slides + notes).

\newcommand{\AIMLRepoURL}{https://github.com/tali7c/Elements-of-AIML}

\newcommand{\AIMLLabSources}{%
Provost and Fawcett, \textit{Data Science for Business}.;%
McKinney, \textit{Python for Data Analysis}.;%
WEKA Documentation (Waikato Environment for Knowledge Analysis), accessed 2026-02-28.;%
Matplotlib and Seaborn documentation, accessed 2026-02-28.%
}

\newcommand{\AIMLUnitISources}{%
Rich and Knight, \textit{Artificial Intelligence}, McGraw Hill, 2017.;%
Russell and Norvig, \textit{Artificial Intelligence: A Modern Approach} (4e), Ch. 1--2 (Introduction, Intelligent Agents).;%
Poole and Mackworth, \textit{Artificial Intelligence: Foundations of Computational Agents} (2e), selected sections.%
}

\newcommand{\AIMLUnitIISources}{%
Russell and Norvig, \textit{Artificial Intelligence: A Modern Approach} (4e), Ch. 7--9 (Logical Agents, First-Order Logic, Inference).;%
Huth and Ryan, \textit{Logic in Computer Science} (2e), selected sections on propositional and first-order logic.;%
Genesereth and Nilsson, \textit{Logical Foundations of Artificial Intelligence}, selected chapters.%
}

\newcommand{\AIMLUnitIIISources}{%
Mueller and Massaron, \textit{Machine Learning for Dummies}, For Dummies, 2016.;%
Mitchell, \textit{Machine Learning}, selected chapters on learning basics and evaluation.;%
Scikit-learn documentation, accessed 2026-02-28.%
}

\newcommand{\AIMLUnitIVSources}{%
Mueller and Massaron, \textit{Machine Learning for Dummies}, For Dummies, 2016.;%
James, Witten, Hastie, and Tibshirani, \textit{An Introduction to Statistical Learning} (2e), selected chapters.;%
Scikit-learn documentation, accessed 2026-02-28.%
}

\newcommand{\AIMLUnitVSources}{%
NIST, \textit{AI Risk Management Framework (AI RMF 1.0)}, 2023.;%
Barocas, Hardt, and Narayanan, \textit{Fairness and Machine Learning} (online book), selected sections.;%
Knaflic, \textit{Storytelling with Data}, selected chapters on communicating results.%
}
%
}{%
  % Faculty copies live under `private/` so their compile CWD is different.
  \IfFileExists{../../../../public/_shared/aiml-sources.tex}{%
    % Source strings for Elements of AIML (used by slides + notes).

\newcommand{\AIMLRepoURL}{https://github.com/tali7c/Elements-of-AIML}

\newcommand{\AIMLLabSources}{%
Provost and Fawcett, \textit{Data Science for Business}.;%
McKinney, \textit{Python for Data Analysis}.;%
WEKA Documentation (Waikato Environment for Knowledge Analysis), accessed 2026-02-28.;%
Matplotlib and Seaborn documentation, accessed 2026-02-28.%
}

\newcommand{\AIMLUnitISources}{%
Rich and Knight, \textit{Artificial Intelligence}, McGraw Hill, 2017.;%
Russell and Norvig, \textit{Artificial Intelligence: A Modern Approach} (4e), Ch. 1--2 (Introduction, Intelligent Agents).;%
Poole and Mackworth, \textit{Artificial Intelligence: Foundations of Computational Agents} (2e), selected sections.%
}

\newcommand{\AIMLUnitIISources}{%
Russell and Norvig, \textit{Artificial Intelligence: A Modern Approach} (4e), Ch. 7--9 (Logical Agents, First-Order Logic, Inference).;%
Huth and Ryan, \textit{Logic in Computer Science} (2e), selected sections on propositional and first-order logic.;%
Genesereth and Nilsson, \textit{Logical Foundations of Artificial Intelligence}, selected chapters.%
}

\newcommand{\AIMLUnitIIISources}{%
Mueller and Massaron, \textit{Machine Learning for Dummies}, For Dummies, 2016.;%
Mitchell, \textit{Machine Learning}, selected chapters on learning basics and evaluation.;%
Scikit-learn documentation, accessed 2026-02-28.%
}

\newcommand{\AIMLUnitIVSources}{%
Mueller and Massaron, \textit{Machine Learning for Dummies}, For Dummies, 2016.;%
James, Witten, Hastie, and Tibshirani, \textit{An Introduction to Statistical Learning} (2e), selected chapters.;%
Scikit-learn documentation, accessed 2026-02-28.%
}

\newcommand{\AIMLUnitVSources}{%
NIST, \textit{AI Risk Management Framework (AI RMF 1.0)}, 2023.;%
Barocas, Hardt, and Narayanan, \textit{Fairness and Machine Learning} (online book), selected sections.;%
Knaflic, \textit{Storytelling with Data}, selected chapters on communicating results.%
}
%
  }{%
    % Source strings for Elements of AIML (used by slides + notes).

\newcommand{\AIMLRepoURL}{https://github.com/tali7c/Elements-of-AIML}

\newcommand{\AIMLLabSources}{%
Provost and Fawcett, \textit{Data Science for Business}.;%
McKinney, \textit{Python for Data Analysis}.;%
WEKA Documentation (Waikato Environment for Knowledge Analysis), accessed 2026-02-28.;%
Matplotlib and Seaborn documentation, accessed 2026-02-28.%
}

\newcommand{\AIMLUnitISources}{%
Rich and Knight, \textit{Artificial Intelligence}, McGraw Hill, 2017.;%
Russell and Norvig, \textit{Artificial Intelligence: A Modern Approach} (4e), Ch. 1--2 (Introduction, Intelligent Agents).;%
Poole and Mackworth, \textit{Artificial Intelligence: Foundations of Computational Agents} (2e), selected sections.%
}

\newcommand{\AIMLUnitIISources}{%
Russell and Norvig, \textit{Artificial Intelligence: A Modern Approach} (4e), Ch. 7--9 (Logical Agents, First-Order Logic, Inference).;%
Huth and Ryan, \textit{Logic in Computer Science} (2e), selected sections on propositional and first-order logic.;%
Genesereth and Nilsson, \textit{Logical Foundations of Artificial Intelligence}, selected chapters.%
}

\newcommand{\AIMLUnitIIISources}{%
Mueller and Massaron, \textit{Machine Learning for Dummies}, For Dummies, 2016.;%
Mitchell, \textit{Machine Learning}, selected chapters on learning basics and evaluation.;%
Scikit-learn documentation, accessed 2026-02-28.%
}

\newcommand{\AIMLUnitIVSources}{%
Mueller and Massaron, \textit{Machine Learning for Dummies}, For Dummies, 2016.;%
James, Witten, Hastie, and Tibshirani, \textit{An Introduction to Statistical Learning} (2e), selected chapters.;%
Scikit-learn documentation, accessed 2026-02-28.%
}

\newcommand{\AIMLUnitVSources}{%
NIST, \textit{AI Risk Management Framework (AI RMF 1.0)}, 2023.;%
Barocas, Hardt, and Narayanan, \textit{Fairness and Machine Learning} (online book), selected sections.;%
Knaflic, \textit{Storytelling with Data}, selected chapters on communicating results.%
}
%
  }%
}


\title{Elements of AIML\\Unit 02 -- Lecture 01 Notes\\Propositional Logic Foundations (Entailment)}
\author{Tofik Ali}
\date{\today}

\begin{document}
\maketitle
\tableofcontents

\section{Lecture Overview}
These notes introduce \textbf{propositional logic} as a tool to represent knowledge and to reason correctly from that knowledge.
In AI, we often store:
\begin{itemize}[leftmargin=*]
  \item \textbf{facts} (what is currently true), and
  \item \textbf{rules} (how truths imply other truths),
\end{itemize}
and then ask queries such as: ``Given my knowledge base, must this query be true?''.

The key concept for that question is \textbf{entailment}:
\[
  KB \models \alpha
  \quad\text{(every model of $KB$ is a model of $\alpha$).}
\]

This lecture also introduces \textbf{CNF} (conjunctive normal form), because many automated reasoning methods expect knowledge in a standard clause-like form (resolution is covered later in this unit).
\notesources{\AIMLUnitIISources}

\section{How to use these notes (self-study)}
Read the sections in order.
When you see a worked example, try to solve it yourself before reading the solution.
For practice problems, write your full reasoning (not only the final answer); in logic, the \emph{justification} is part of the learning outcome.

\section{Learning Outcomes}
\begin{itemize}[leftmargin=*]
  \item Define propositions and logical connectives; compute truth values using truth tables.
  \item Classify formulas as valid (tautology), satisfiable, or unsatisfiable.
  \item Explain semantic entailment $KB \models \alpha$ and test it for small KBs using models/truth tables.
  \item Convert simple propositional formulas to CNF and explain why CNF is useful.
  \item Explain (at a high level) why propositional logic is insufficient for many-object relational domains.
\end{itemize}

\section{Keywords}
\begin{itemize}[leftmargin=*]
  \item \textbf{Proposition:} a declarative statement that is either True or False.
  \item \textbf{WFF (well-formed formula):} a syntactically valid logic expression.
  \item \textbf{Model:} a truth assignment that makes a sentence true.
  \item \textbf{Entailment ($\models$):} truth in all models (semantic consequence).
  \item \textbf{Validity:} true in all models (tautology).
  \item \textbf{Satisfiable:} true in at least one model.
  \item \textbf{CNF:} conjunction (AND) of disjunctions (OR-clauses).
\end{itemize}

\section{Abbreviations and symbols}
\begin{itemize}[leftmargin=*]
  \item $KB$: knowledge base (set of sentences/facts/rules).
  \item $\alpha$: query/goal sentence.
  \item $\lnot,\land,\lor,\rightarrow,\leftrightarrow$: NOT, AND, OR, IMPLIES, IFF.
  \item $\models$: entailment (semantic); $\vdash$: derivation/provability (syntactic).
  \item \textbf{CNF}: conjunctive normal form; \textbf{literal}: $p$ or $\lnot p$; \textbf{clause}: OR of literals.
\end{itemize}

\section{Why logic in AI? (KBs and queries)}
\subsection{Knowledge base and query}
A \textbf{knowledge base (KB)} is a set of sentences describing a world.
For example:
\begin{itemize}[leftmargin=*]
  \item Fact: $p$ = ``It is raining.''
  \item Rule: $p \rightarrow q$ = ``If it is raining, then the ground is wet.''
\end{itemize}
A \textbf{query} $\alpha$ is what we want to know, e.g., $q$ = ``The ground is wet.''.

\subsection{Syntax vs semantics (the important separation)}
\begin{itemize}[leftmargin=*]
  \item \textbf{Syntax}: which strings count as legal formulas (grammar).
  \item \textbf{Semantics}: what those formulas mean (truth in models).
\end{itemize}
We can manipulate syntax mechanically (proofs) and still preserve meaning if our inference rules are sound.

\section{Propositional logic}
\subsection{Syntax: propositions and well-formed formulas}
Propositional logic uses propositional symbols (atoms) like $p,q,r$.
We build complex formulas using connectives:
\[
  \lnot p,\; (p\land q),\; (p\lor q),\; (p\rightarrow q),\; (p\leftrightarrow q)
\]
A formula is a \textbf{well-formed formula (WFF)} if it follows the grammar (symbols, connectives, parentheses).

\subsection{Semantics: truth assignments and models}
A \textbf{truth assignment} $v$ maps each propositional symbol to $T$ or $F$.
The truth of a complex formula is computed using the connective truth tables.

If a formula $\varphi$ is true under assignment $v$, we say $v$ is a \textbf{model} of $\varphi$ and write $v \models \varphi$.
If $v$ makes all sentences in $KB$ true, then $v$ is a model of $KB$ (treating $KB$ as a conjunction of its sentences).

\subsection{Truth tables and classification of formulas}
Truth tables list all possible assignments and compute the formula value.
With $n$ propositional symbols, there are $2^n$ rows.

Using a truth table we can classify:
\begin{itemize}[leftmargin=*]
  \item \textbf{Tautology/valid}: always true (e.g., $p\lor \lnot p$).
  \item \textbf{Contradiction/unsatisfiable}: always false (e.g., $p\land \lnot p$).
  \item \textbf{Satisfiable}: true in at least one assignment (has a model).
  \item \textbf{Unsatisfiable KB}: no assignment satisfies all KB sentences (KB is inconsistent).
\end{itemize}

\subsection{Implication ($p\rightarrow q$) and common confusions}
In everyday language, ``if'' sometimes implies causation.
In propositional logic, implication is purely truth-functional:
\[
  (p\rightarrow q) \text{ is false only when $p$ is true and $q$ is false.}
\]
The most useful equivalence:
\[
  (p\rightarrow q) \equiv (\lnot p \lor q)
\]
This equivalence is used in CNF conversion and SAT/resolution methods.

\section{Entailment ($KB \models \alpha$)}
\subsection{Definition (semantic)}
$KB$ entails $\alpha$ (written $KB \models \alpha$) iff:
\begin{quote}
  In \emph{every} model where $KB$ is true, $\alpha$ is also true.
\end{quote}
Equivalently: there is \textbf{no counterexample model} where $KB$ is true and $\alpha$ is false.

\subsection{How to test entailment in propositional logic}
For propositional logic, the following are equivalent (and can be checked by truth tables for small KBs):
\begin{itemize}[leftmargin=*]
  \item $KB \models \alpha$
  \item $(KB \rightarrow \alpha)$ is valid (a tautology), where $KB$ is treated as a conjunction of its sentences
  \item $KB \land \lnot \alpha$ is \textbf{unsatisfiable}
\end{itemize}
The last form is important because many automated reasoning methods reduce inference to satisfiability.

\subsection{Entailment vs implication vs derivation (do not mix them up)}
\begin{itemize}[leftmargin=*]
  \item \textbf{Implication} ($\rightarrow$) is a connective inside the language.
  \item \textbf{Entailment} ($\models$) is a semantic relationship (about all models).
  \item \textbf{Derivation} ($\vdash$) is a syntactic relationship (there exists a proof using rules).
\end{itemize}

\paragraph{Soundness and completeness (intuition).}
A proof system is \textbf{sound} if it never proves a sentence that is not entailed.
It is \textbf{complete} if it can prove every sentence that is entailed.

\subsection{Worked example: model-based reasoning}
Let:
\[
  KB=\{(p\rightarrow q),\; (q\rightarrow r),\; p\},\qquad \alpha=r
\]
Reasoning using rules:
\begin{enumerate}[leftmargin=*]
  \item From $(p\rightarrow q)$ and $p$, derive $q$ (modus ponens).
  \item From $(q\rightarrow r)$ and $q$, derive $r$.
\end{enumerate}
So every model of $KB$ must make $r$ true, hence $KB \models r$.

If you wanted to check via models (truth table), you would enumerate assignments of $(p,q,r)$, keep those that satisfy all three KB sentences, and verify that in all such rows, $r=T$.

\section{Normal forms and CNF (why we care)}
\subsection{Literals, clauses, and CNF}
\begin{itemize}[leftmargin=*]
  \item A \textbf{literal} is $p$ or $\lnot p$.
  \item A \textbf{clause} is an OR of literals, e.g., $(\lnot p \lor q \lor r)$.
  \item \textbf{CNF} is an AND of clauses, e.g.,
  \[
    (\lnot p \lor q)\land (p \lor r)\land (\lnot r)
  \]
\end{itemize}

\subsection{CNF conversion recipe (propositional)}
To convert a formula to CNF:
\begin{enumerate}[leftmargin=*]
  \item Eliminate $\leftrightarrow$ and $\rightarrow$.
  \item Push negations inward (De Morgan + double negation).
  \item Distribute $\lor$ over $\land$ until you get an AND of ORs.
  \item Flatten nested AND/OR groups.
\end{enumerate}

\subsection{Worked CNF example (distribution is required)}
Convert $p \rightarrow (q \land r)$:
\begin{enumerate}[leftmargin=*]
  \item Eliminate implication:
  \[
    \lnot p \lor (q \land r)
  \]
  \item Distribute $\lor$ over $\land$:
  \[
    (\lnot p \lor q) \land (\lnot p \lor r)
  \]
\end{enumerate}
This is now in CNF (two clauses).

\section{Limits of propositional logic (preview of predicate logic)}
Propositional logic is not designed for ``many objects + relations'' worlds.
It lacks:
\begin{itemize}[leftmargin=*]
  \item \textbf{relational structure} (e.g., $Parent(x,y)$),
  \item \textbf{quantifiers} (``all'' / ``some''),
  \item \textbf{compact general rules} that apply to arbitrarily many objects.
\end{itemize}
Predicate (first-order) logic addresses this by adding objects, predicates, and quantifiers.
We cover that fully in \textbf{Unit~02 Lecture~02}.

\section{Common pitfalls (and how to avoid them)}
\begin{itemize}[leftmargin=*]
  \item \textbf{Mixing up $\models$ and $\rightarrow$:} $KB \models \alpha$ is a semantic relationship (about all models); $KB \rightarrow \alpha$ is a formula inside the language.
  \item \textbf{Forgetting the meaning of ``model of a KB'':} a model of $KB$ is a truth assignment that makes \emph{every} sentence in $KB$ true (treat $KB$ like a conjunction).
  \item \textbf{Treating ``satisfiable'' as ``true'':} satisfiable means ``true in at least one assignment''; it does not mean ``true in the real world''.
  \item \textbf{CNF conversion mistakes:} eliminate $\leftrightarrow$ and $\rightarrow$ first, then push negations inward, \emph{then} distribute $\lor$ over $\land$.
  \item \textbf{Implication intuition trap:} $p\rightarrow q$ does not mean causation; it is false only when $p$ is true and $q$ is false.
\end{itemize}

\section{Practice (with answers)}
\subsection*{P1. Build a truth table and decide if $(p\rightarrow q)\leftrightarrow(\lnot p \lor q)$ is valid.}
\textbf{Answer:} It is valid. Both expressions have identical truth values for all assignments of $(p,q)$.

\subsection*{P2. Give a satisfying assignment for $(p \lor q)\land (\lnot p \lor r)\land (\lnot r)$.}
\textbf{Answer:} $r=F$ (from $\lnot r$), then $\lnot p$ must be true so $p=F$, then $q=T$ to satisfy $(p\lor q)$. One model: $(p=F,q=T,r=F)$.

\subsection*{P3. Show that $KB=\{(p\rightarrow q),(q\rightarrow r),p\}$ entails $r$.}
\textbf{Answer:} Modus ponens yields $q$, then yields $r$, so $KB\models r$.

\subsection*{P4. Convert to CNF: $p \leftrightarrow q$.}
\textbf{Answer:} $(p\leftrightarrow q)\equiv (p\rightarrow q)\land(q\rightarrow p)\equiv (\lnot p\lor q)\land(\lnot q\lor p)$.

\subsection*{P5. Convert to CNF: $p \rightarrow (q \land r)$.}
\textbf{Answer:} $(\lnot p\lor q)\land(\lnot p\lor r)$.

\subsection*{P6. Give one limitation of propositional logic in real knowledge bases.}
\textbf{Answer:} it cannot express general rules like ``all students...'' without writing separate symbols/rules for each object.

\section{Exit questions (with answers)}
\subsection*{Q1. What does $KB \models \alpha$ mean (in one sentence)?}
\textbf{Answer:} In every model where all sentences of $KB$ are true, $\alpha$ is also true.

\subsection*{Q2. For a small KB, how can we test entailment using models/truth tables?}
\textbf{Answer:} Enumerate truth assignments; keep those that satisfy $KB$; check that $\alpha$ is true in all of them (no counterexample).

\subsection*{Q3. What is one common confusion about $p \rightarrow q$?}
\textbf{Answer:} Confusing it with causation; logically it is false only when $p$ is true and $q$ is false.

\subsection*{Q4. Give one limitation of propositional logic that motivates predicate logic.}
\textbf{Answer:} Propositional logic cannot compactly express relations and quantifiers (many-object rules).

\section{Summary}
Propositional logic provides a precise language for facts and rules, and \textbf{entailment} defines what it means for a query to follow from a knowledge base.
Truth tables provide a brute-force entailment test for small problems.
CNF is a standard representation that prepares formulas for automated inference procedures.
\notesources{\AIMLUnitIISources}

\section{Further Reading}
\begin{itemize}[leftmargin=*]
  \item Russell and Norvig, \textit{Artificial Intelligence: A Modern Approach} (4e), chapters on logical agents.
  \item Huth and Ryan, \textit{Logic in Computer Science} (2e), propositional logic chapters.
\end{itemize}
\notesources{\AIMLUnitIISources}

\section{Next Lecture Preview}
\textbf{Next (Unit~02 Lecture~02):} Predicate (first-order) logic --- predicates, quantifiers, translating English to logic, and scope pitfalls.

\end{document}
